\section{Introduction}

Large Language Models (LLMs) such as GPT-4, Gemini, and Claude have evolved from narrow text-completion tools into broadly capable reasoning engines, surpassing human performance on tasks ranging from legal exams~\cite{openai2023gpt4} to software engineering challenges~\cite{swebench2023}. However, most of these benchmarks are \textit{single-turn} interactions. Real-world deployments, in contrast, demand persistence and adaptability: evolving software repositories over weeks, synthesizing iterative scientific reviews, or mediating multistakeholder corporate decisions. These are inherently multiturn, long-horizon, and collaborative, pushing beyond what any single monolithic model, no matter how large, can reliably achieve.

Multi-agent LLM systems have emerged as a promising direction to scale intelligence by composing the complementary skills of multiple agents. For example, one agent might specialize in writing unit tests, another in literature surveys, and a third in policy analysis. By enabling interaction between such experts, these systems can address tasks outside the scope of any individual model. Yet, their deployment remains hampered by three persistent bottlenecks:

\textbf{(1) Dynamic task routing.} Existing orchestrators often rely on fixed pipelines or static vector similarity rules. These approaches are brittle: When task requirements change or agent performance degrades, the system continues to misroute, invoking wrong experts, and compounding errors.

\textbf{(2) Credit assignment across long horizons.} In extended dialogues, failures may not be visible until dozens of turns later. Without timely and fine-grained feedback, the system cannot effectively demote underperforming agents or up-weight reliable ones, leading to stagnation in routing efficiency.

\textbf{(3) Cold-start inefficiency.} When a new task arrives, the system has no prior evidence about which agents are competent. Without mechanisms to transfer knowledge from previous interactions, early routing is effectively random, wasting tokens and agent calls before the system can differentiate agent quality.

We propose \method, a lightweight and training-free controller for multi-agent LLM systems.
\method\ wraps any pool of agents with four components:
(1) belief-guided delegation via Thompson sampling to prioritize agents with historically positive marginal contributions;
(2) calibrated reflection via a judge that triggers credit assignment and re-routing;
(3) evidence-checked selection (rather than averaging); and
(4) memory-aware priors to reduce cold-start inefficiency and context bloat.

By combining these elements into a recursive loop, \method\ enables efficient, belief-guided routing across a pool of heterogeneous agents. In controlled experiments on split-knowledge tasks, we show that belief-guided routing reduces token usage by 28\%, agent calls by 17\%, and time-to-success by 19\% compared to an otherwise identical system with uniform random selection, while maintaining matched task success rates.

\paragraph{Contributions.}
We make the following contributions:
\begin{itemize}
    \item We introduce \textsc{REDEREF}, a training-free controller for multi-agent LLM systems that
    improves routing efficiency under recursive retry without requiring fine-tuning or centralized training.
    \item We reinterpret belief updates as modeling an agent's probability of providing a
    \emph{positive marginal contribution}, resolving credit assignment in compositional multi-agent tasks.
    \item We demonstrate empirically that belief-guided delegation substantially reduces token usage,
    agent calls, and time-to-success compared to random recursive delegation at matched success rates.
    \item We analyze robustness to agent impairment and judge miscalibration, showing graceful
    degradation rather than catastrophic routing failures.
\end{itemize}
